\subsection{基的变换}

命题4. 设 E 是一个右 A-模，具有有限基 \((e_i)_{1 \leqslant i \leqslant n}\) ，由 \(n\) 个元素组成。

要使由 \(n\) 个元素 \(e_i' = \sum_{j=1} e_j a_{ji} (1 \leqslant i \leqslant n)\) 组成的族构成 \(\mathbf{E}\) 的基，必要且充分的条件是方阵 \(P = (a_{ji})\) （\(n\) 阶）为可逆。

P 恰好是 E 的自同态 u 关于基 \((e_{i})\) 的矩阵，其中 u 由 \(u(e_{i}) = e_{i}^{\prime} (1 \leqslant i \leqslant n)\) 定义。现在，要使 u 成为 E 的自同构，必要且充分的是 \((u(e_{i}))\) 构成 E 的基（§ 1，第 11 号，命题 17 的推论 3）；由此得命题。

可逆矩阵 P 称为从基 \((e_{i})\) 到基 \((e_{i}^{\prime})\) 的过渡矩阵。它也可以解释为恒等映射 \(1_{E}\) 关于基 \((e_{i}^{\prime})\) 和 \((e_{i})\) （按此顺序）的矩阵；于是显然，从基 \((e_{i}^{\prime})\) 到基 \((e_{i})\) 的过渡矩阵是 P 的逆 \(P^{-1}\) 。

命题5. 设 \((e_i)\) ， \((e_i')\) 是由 \(n\) 个元素构成的 \(\mathbf{E}\) 的两个基， \(P\) 是从 \((e_i)\) 到 \((e_i')\) 的过渡矩阵。若 \((e_i^*)\) 和 \((e_i'^*)\) 分别是 \((e_i)\) 和 \((e_i')\) 的对偶基，则从 \((e_i^*)\) 到 \((e_i'^*)\) 的过渡矩阵是逆步矩阵 \(^t P^{-1}\) （即 \(P\) 的逆步）。

恒等映射 \(1_{E}\) 的转置是恒等映射 \(1_{E^{*}}\) ；根据第 4 号命题 3， \(1_{E^{*}}\) 关于基 \((e_{i}^{\prime*})\) 和 \((e_{i}^{*})\) （按此顺序）的矩阵是 \(1_{E}\) 关于基 \((e_{i})\) 和 \((e_{i}^{\prime})\) （按此顺序）的矩阵的转置，即 \(P^{-1}\) 的转置。

命题6. 设 E 和 F 是两个右 A-模， \((e_i)\) 和 \((e_i')\) 是 E 的含 \(n\) 个元素的两个基， \((f_j)\) 和 \((f_j')\) 是 \(F\) 的含 \(m\) 个元素的两个基， \(P\) 是从 \((e_i)\) 到 \((e_i')\) 的过渡矩阵， \(Q\) 是从 \((f_j)\) 到 \((f_j')\) 的过渡矩阵。对于每个线性映射 \(u\) （由 \(E\) 到 \(F\) ），设 \(M(u)\) 是 \(u\) 关于基 \((e_i)\) 和 \((f_j)\) 的矩阵， \(M'(u)\) 是 \(u\) 关于基 \((e_i')\) 和 \((f_i')\) 的矩阵；则

\[
M ^ {\prime} (u) = Q ^ {- 1} M (u) P.\tag{33 D}
\]

我们可以写成 \(u = 1_{\mathbb{F}} \circ u \circ 1_{\mathbb{E}}\) 。当 \(1_{\mathbb{E}}\) 的矩阵取为关于 \((e_i')\) 和 \((e_i)\) 、 \(u\) 的矩阵取为关于 \((e_i)\) 和 \((f_i)\) 、而 \(1_{\mathbb{F}}\) 的矩阵取为关于 \((f_i)\) 和 \((f_j')\) 时，公式 (33) 立即由第 4 号命题 2 的推论得出。

推论1. 若 \(u\) 是 \(E\) 的自同态， \(M(u)\) 和 \(M'(u)\) 是它分别关于基 \((e_1)\) 和 \((e_1')\) 的矩阵，则

\[
M ^ {\prime} (u) = P ^ {- 1} M (u) P.\tag{34 D}
\]

推论2. 如果 \(M(x)\) 和 \(M'(x)\) 是同一元素 \(x \in E\) 分别关于基 \((e_i)\) 和 \((e_i')\) 的单列矩阵，则

\[
M (x) = P. M ^ {\prime} (x).\tag{35 D}
\]

这是命题 6 应用于映射 \(\theta_{x}:a\mapsto xa\) （由 \(\mathbf{A}_d\) 到 \(\mathbf{E}\) ）的一个特例（第 4 号，注记 1）。

公式 (35) 等价于

\[
x _ {i} = \sum_ {j = 1} ^ {n} a _ {i j} x _ {j} ^ {\prime} (1 \leqslant i \leqslant n)\tag{36 D}
\]

对于元素 \(x_{i}\) 和 \(x_{i}^{\prime}\) （分别属于矩阵 \(M(x)\) 和 \(M^{\prime}(x)\) ）。公式 (36 D) 称为坐标变换公式。注意，它们把 x 相对于“旧”基 ( \(e_{i}\) ) 的分量表示成 x 相对于“新”基 ( \(e_{i}^{\prime}\) ) 的分量以及 P 的元素（即“新”基相对于“旧”基的分量）的函数。

注记。(1) 现在我们从一个左 A-模 E 出发，它有两个基 \((e_{i})\) , \((e_{i}^{\prime})\) ，每个基各有 n 个元素；如果记 \(e_{i}^{\prime} = \sum_{j=1}^{n} a_{ji} e_{i}\) ，则 \(P = (a_{ji})\) 也称为从 \((e_{i})\) 到 \((e_{i}^{\prime})\) 的过渡矩阵；它也是 E 的使 \(u(e_{i}) = e_{i}^{\prime}\) 的自同构 u 关于基 \((e_{i})\) 的矩阵，并且也是 \(1_{E}\) 关于基 \((e_{i}^{\prime})\) 和 \((e_{i})\) （按此顺序）的矩阵。上述结果仍然成立，只需作如下修改：公式(33 D)至(36 D)分别替换为

\[
{ } ^ { t } M ^ { \prime } ( u ) = { } ^ { t } P . { } ^ { t } M ( u ) . { } ^ { t } Q ^ { - 1 }\tag{33 G}
\]

\[
{ } ^ { t } M ^ { \prime } ( u ) = { } ^ { t } P . { } ^ { t } M ( u ) . { } ^ { t } P ^ { - 1 }\tag{34 G}
\]

\[
{ } ^ { t } M ( u ) = { } ^ { t } M ^ { \prime } ( u ) . { } ^ { t } P .\tag{35 G}
\]

\[
\kappa_ {i} = \sum_ {j = 1} ^ {n} x _ {j} ^ {\prime} a _ {i j} \quad (1 \leqslant i \leqslant n).\tag{36 G}
\]

\begin{enumerate}
\def\labelenumi{(\arabic{enumi})}
\setcounter{enumi}{1}
\tightlist
\item
  在命题 4 的假设下，考虑一个元素 \(x^{*} \in \mathbf{E}^{*}\) ；由于从 \((e_i^*)\) 到 \((e_i'^*)\) 的过渡矩阵是 \(^tP^{-1}\) （命题 5），对于矩阵 \(M(x^{*})\) 和 \(M'(x^{*})\) （即 \(x^{*}\) 分别关于这两个基的矩阵），
\end{enumerate}

\[
{ } ^ { t } M ( x ^ { * } ) = { } ^ { t } M ^ { \prime } ( x ^ { * } ) . P ^ { - 1 }
\]

或等价地

\[
{ } ^ { t } M ^ { \prime } ( x ^ { * } ) = { } ^ { t } M ( x ^ { * } ) . P\tag{37 D}
\]

这等价于方程组

\[
x _ {i} ^ {\prime *} = \sum_ {j = 1} ^ {n} x _ {j} ^ {*} a _ {j i} (1 \leqslant i \leqslant n)\tag{38 D}
\]

对于元素 \((x_{i}^{*})\) 和 \((x_{i}^{\prime *})\) （分别属于矩阵 \(M(x^{*})\) 和 \(M'(x^{*})\) ）。对于左 A-模 E，相应的公式为

\[
M ^ {\prime} (x ^ {*}) = ^ {t} P. M (x ^ {*})\tag{37 G}
\]

\[
x _ {i} ^ {\prime *} = \sum_ {j = 1} ^ {n} a _ {j i} x _ {j} ^ {*} (1 \leqslant i \leqslant n).\tag{38 G}
\]

\begin{enumerate}
\def\labelenumi{(\arabic{enumi})}
\setcounter{enumi}{2}
\tightlist
\item
  设 A、B 为两个环， \(\sigma: \mathbf{A} \to \mathbf{B}\) 是从 A 到 B 的一个同态，E 是一个右（相应地，左）A-模， \((e_i)\) , \((e_i')\) 是 E 的两个基，各含 \(n\) 个元素，F 是一个右（相应地，左）B-模， \((f_j)\) , \((f_j')\) 是 F 的两个基，各含 \(m\) 个元素，以及 \(P\) （相应地 \(Q\) ）是从 \((e_i)\) 到 \((e_i')\) （相应地，从 \((f_j)\) 到 \((f_j')\) ）的过渡矩阵。
\end{enumerate}

对于每个半线性映射 \(u: \mathbf{E} \to \mathbf{F}\) （相对于 \(\sigma\) ），设 \(M(u)\) 为 \(u\) 关于 \((e_i)\) 和 \((f_j)\) 的矩阵， \(M'(u)\) 为其关于 \((e_i')\) 和 \((f_j')\) 的矩阵。则

\[
M ^ {\prime} (u) = Q ^ {- 1} M (u) \sigma (P)\tag{39 D}
\]

（相应地

\[
{ } ^ { t } M ^ { \prime } ( u ) = \sigma ( { } ^ { t } P ) . { } ^ { t } M ( u ) . { } ^ { t } Q ^ { - 1 } ) .\tag{39 G}
\]

证明与 (33 D) 和 (33 G) 的证明相同，这次使用公式 (27 D) 和 (27 G)（第 6 号）。

